% =====================================================================
\chapter{Soluciones de los ejercicios}
\label{ap:soluciones}
% =====================================================================

Cada solución se presenta en cinco partes: datos, planteamiento, desarrollo, resultado y comprobación. En los pasos intermedios se redondea a cuatro cifras significativas salvo que se indique otra cosa. Los resultados numéricos de los casos de estudio se han comprobado con el cuaderno \texttt{M24-practica.ipynb}.

% ---------------------------------------------------------------------
\section{Tema 1}

\begin{solucion}{eje:1-elementos}
\datos invernadero de 1~ha con control horario de ventanas, calefacción y CO$_2$; predicción de radiación y temperatura exterior.
\planteamiento clasificar cada elemento según la \cref{def:elementos}.
\pasos
\begin{enumerate}[label=\arabic*.]
  \item Variables de decisión, por hora $t$: apertura de ventanas $v_t\in[0,1]$, potencia de calefacción $u_t\ge 0$ (kW) y caudal de CO$_2$ $g_t\ge 0$ (kg/h).
  \item Parámetros: radiación $R_t$ y temperatura exterior $T^{\mathrm{ext}}_t$ (predichos); precio de la energía y del CO$_2$ (conocidos o predichos); coeficientes del modelo térmico y del modelo de crecimiento (estimados con datos).
  \item Objetivo: $\max\sum_t\bigl[\text{crecimiento estimado}(T_t,\mathrm{CO}_{2,t},R_t)-p^{\mathrm{e}}_tu_t-p^{\mathrm{c}}g_t\bigr]$.
  \item Restricciones duras: $\mathrm{CO}_{2,t}\le 1000$~ppm (seguridad de las personas), límites físicos de los actuadores y la dinámica térmica y de CO$_2$ del invernadero.
  \item Restricción que conviene tratar como blanda: $16\le T_t\le 28$~ºC, con penalización por grado fuera de rango, porque una helada imprevista podría hacerla infactible; si el cultivo sufre daños irreversibles fuera del rango, debe ser dura.
  \item KPIs distintos del objetivo: horas fuera del rango de temperatura, kWh por kilogramo producido, error de la predicción de temperatura interior, número de intervenciones manuales.
\end{enumerate}
\resultado \textbf{decisiones $v_t,u_t,g_t$; parámetros predichos $R_t,T^{\mathrm{ext}}_t$; objetivo crecimiento menos coste; CO$_2$ como restricción dura y temperatura como blanda; KPIs de confort, eficiencia, predicción y supervisión.}
\comprobacion cada elemento es o bien elegible por el sistema (decisión) o bien externo (parámetro); ningún KPI coincide con el objetivo.
\end{solucion}

\begin{solucion}{eje:1-mapa}
\datos seis problemas.
\planteamiento identificar en cada uno el tipo de variables, la linealidad y el coste de evaluación (\cref{tab:mapa,fig:arbol}).
\pasos
\begin{enumerate}[label=\arabic*.]
  \item (a) Variables continuas, todo lineal: \emph{PL}.
  \item (b) Decisión binaria de apertura con costes lineales: \emph{PLEM} (localización con costes fijos).
  \item (c) Caja negra de 2~h por evaluación y 8 variables continuas: \emph{optimización bayesiana}.
  \item (d) Secuenciación con precedencias y recursos disyuntivos: \emph{programación por restricciones} (CP-SAT).
  \item (e) Objetivo cuadrático convexo y restricciones lineales: \emph{optimización convexa} (programa cuadrático).
  \item (f) Decisiones secuenciales con realimentación cada minuto: \emph{control predictivo o aprendizaje por refuerzo}.
\end{enumerate}
\resultado \textbf{(a) PL; (b) PLEM; (c) optimización bayesiana; (d) CP; (e) convexa; (f) control o RL.}
\comprobacion en ningún caso se usa una metaheurística genérica donde existe un método que explota la estructura.
\end{solucion}

\begin{solucion}{eje:1-abs}
\datos $\min\abs{x_1}+2\abs{x_2}$ s.a. $x_1+x_2\ge 4$, $x_1-x_2\le 2$.
\planteamiento \cref{prop:reform} (iii): $x_j=x_j^{+}-x_j^{-}$; en el óptimo $\abs{x_j}=x_j^{+}+x_j^{-}$.
\pasos
\begin{enumerate}[label=\arabic*.]
  \item PL: $\min x_1^{+}+x_1^{-}+2x_2^{+}+2x_2^{-}$ s.a. $x_1^{+}-x_1^{-}+x_2^{+}-x_2^{-}\ge 4$, $x_1^{+}-x_1^{-}-x_2^{+}+x_2^{-}\le 2$, todas $\ge 0$.
  \item Es mejor hacer grande $x_1$ (coste 1) que $x_2$ (coste 2), pero $x_1\le 2+x_2$.
  \item Con $x_1,x_2\ge 0$, el objetivo es $x_1+2x_2$; sobre la recta $x_1=2+x_2$, la restricción $x_1+x_2\ge 4$ da $x_2\ge 1$; el mínimo es $x_2=1$, $x_1=3$, con valor $3+2=5$.
  \item Con $x_2<0$ haría falta $x_1\ge 4-x_2>4$ y $x_1\le 2+x_2<2$: imposible. Con $x_1<0$ haría falta $x_2>4$ y coste $>8$.
\end{enumerate}
\resultado $\bm{x^{*}=(3,1)}$, valor óptimo \textbf{5}.
\comprobacion $3+1=4\ge 4$ y $3-1=2\le 2$ (ambas activas); HiGHS da el mismo resultado sobre el PL del paso~1.
\end{solucion}

\begin{solucion}{eje:1-correos}
\datos demandas $r=(17,13,15,19,14,16,11)$ de lunes a domingo; cinco días seguidos de trabajo.
\planteamiento variables $x_j$, empleados que empiezan el día $j$; cada día lo cubren los que empezaron ese día o en los cuatro anteriores.
\pasos
\begin{enumerate}[label=\arabic*.]
  \item (a) $\min\sum_{j=1}^{7}x_j$ s.a. $x_1+x_4+x_5+x_6+x_7\ge 17$ (lunes), $x_1+x_2+x_5+x_6+x_7\ge 13$, \dots, $x_3+x_4+x_5+x_6+x_7\ge 11$ (domingo), $x_j\ge 0$ enteros. Con la variable «empleados presentes el día $j$» no se puede expresar la regla de cinco días seguidos; con la variable de inicio la regla queda incorporada en los coeficientes.
  \item (b) Cada $x_j$ aparece en exactamente cinco restricciones. Sumándolas: $5\sum_jx_j\ge\sum_jr_j=105$, luego $\sum_jx_j\ge 21$.
  \item (c) La relajación ($22{,}33$) es una cota mejor que 21 porque usa las restricciones por separado; como el número de empleados es entero, el óptimo no puede bajar de $\lceil 22{,}33\rceil=23$; al encontrar una solución entera con 23, la cota redondeada demuestra que es óptima.
\end{enumerate}
\resultado \textbf{cota por suma: 21; relajación: $22{,}33$; óptimo entero: 23}, certificado por $\lceil 22{,}33\rceil=23$.
\comprobacion la solución entera $x=(7,5,1,8,0,2,0)$ suma 23 y cubre el lunes con $7+8+0+2+0=17$ y el jueves con $7+5+1+8+0=21\ge 19$.
\end{solucion}

\begin{solucion}{eje:1-piensos}
\datos vértices $(0,0)$, $(150,0)$, $(100,100)$, $(40,160)$, $(0,160)$; nuevo margen $40E+45P$.
\planteamiento evaluar el objetivo en los vértices (\cref{teo:vertice}) y estudiar el cono de las restricciones activas.
\pasos
\begin{enumerate}[label=\arabic*.]
  \item Valores: $0$; $6000$; $4000+4500=8500$; $1600+7200=8800$; $7200$.
  \item El óptimo es $(40,160)$ con 8800~€. KPIs: molino $0{,}5\cdot 200=100$~h (100\,\%); granuladora $24+48=72$~h (80\,\%); demanda de puesta cubierta al 100\,\%.
  \item $(100,100)$ es óptimo si el gradiente $(40,m_P)$ es combinación no negativa de las normales activas $(1,1)$ y $(2,1)$: $(40,m_P)=\alpha(1,1)+\beta(2,1)$, es decir, $\beta=40-m_P\ge 0$ y $\alpha=2m_P-40\ge 0$.
\end{enumerate}
\resultado \textbf{con $m_P=45$: $(E,P)=(40,160)$ y 8800~€/semana; $(100,100)$ es óptimo para $20\le m_P\le 40$~€/t.}
\comprobacion con $m_P=40$ empatan $(100,100)$ y $(40,160)$: $4000+4000=8000=1600+6400$; con $m_P=20$ empatan $(100,100)$ y $(150,0)$: $6000$ en ambos.
\end{solucion}

\begin{solucion}{eje:1-minimax}
\datos $f_1=x-1$, $f_2=-2x+8$, $f_3=0{,}5x+1$, $x\in[0,10]$.
\planteamiento \cref{prop:minimax}: $\min t$ s.a. $f_i(x)\le t$, $0\le x\le 10$.
\pasos
\begin{enumerate}[label=\arabic*.]
  \item PL: $\min t$ s.a. $x-1\le t$, $-2x+8\le t$, $0{,}5x+1\le t$, $0\le x\le 10$.
  \item La envolvente superior decrece mientras domina $f_2$ y crece después. $f_2=f_3$ en $-2x+8=0{,}5x+1$, es decir, $x=2{,}8$, con valor $2{,}4$.
  \item En $x=2{,}8$, $f_1=1{,}8<2{,}4$: la envolvente es $\max(f_2,f_3)$ en ese punto.
\end{enumerate}
\resultado $\bm{x^{*}=2{,}8}$, pérdida máxima \textbf{2,4~kW}.
\comprobacion en $x=2{,}7$, $\max=f_2=2{,}6$; en $x=2{,}9$, $\max=f_3=2{,}45$; ambos mayores que $2{,}4$.
\end{solucion}

\begin{solucion}{eje:1-desviaciones}
\datos $\sum_ja_{ij}x_j+d_i-e_i=r_i$, $d_i,e_i\ge 0$; objetivo $\sum_i(\alpha d_i+\beta e_i)$ con $\alpha,\beta>0$.
\planteamiento por reducción al absurdo: si $d_i,e_i>0$, se construye una solución mejor.
\pasos
\begin{enumerate}[label=\arabic*.]
  \item Sea una solución óptima con $d_i>0$ y $e_i>0$. Sea $\delta=\min(d_i,e_i)>0$.
  \item Se define $d_i'=d_i-\delta\ge 0$, $e_i'=e_i-\delta\ge 0$; la restricción se mantiene porque $d_i'-e_i'=d_i-e_i$.
  \item El objetivo baja en $(\alpha+\beta)\delta>0$, en contra de la optimalidad. Luego $d_ie_i=0$.
  \item Con $5d_i^{2}+e_i^{2}$ el argumento es el mismo: $5(d_i-\delta)^{2}+(e_i-\delta)^{2}<5d_i^{2}+e_i^{2}$ porque ambos términos decrecen al reducir sus bases positivas.
\end{enumerate}
\resultado \textbf{en el óptimo $d_ie_i=0$ para todo $i$, con costes lineales positivos y también con el coste cuadrático.}
\comprobacion es el argumento de la clase~2 de 15.053: con $d_1=1$ y $e_1=4$, restar 1 a ambos da $d_1=0$, $e_1=3$, mejor.
\end{solucion}

\begin{solucion}{eje:1-diseno}
\datos invernadero del \cref{eje:1-elementos}; modelo térmico lineal en los controles.
\planteamiento recorrer los bloques de la \cref{fig:anatomia}.
\pasos
\begin{enumerate}[label=\arabic*.]
  \item \emph{Datos}: sensores interiores (temperatura, humedad, CO$_2$, radiación), estación exterior, consumo de energía, estado de actuadores y previsión meteorológica, con marca de tiempo y control de calidad.
  \item \emph{Predicción}: modelo de regresión o red pequeña para radiación y temperatura exterior de las próximas 24~h a partir de la previsión oficial (corrige su sesgo local); modelo identificado del invernadero (lineal en los controles) validado en semanas no usadas para ajustarlo.
  \item \emph{Optimización}: control predictivo con horizonte de 24~h repetido cada hora; con modelo lineal y coste lineal o cuadrático convexo es un PL o un problema cuadrático, resoluble en milisegundos con HiGHS u OSQP.
  \item \emph{Salvaguardas}: un controlador de respaldo por reglas que abre ventanas si el CO$_2$ supera 1000~ppm o la temperatura 30~ºC, independiente del optimizador; límites de variación de los actuadores.
  \item \emph{Monitorización}: error de la predicción de temperatura interior, horas fuera de rango, energía por kilogramo, frecuencia de actuación de las salvaguardas y tiempo de cálculo; si el error crece de forma sostenida, se reidentifica el modelo.
\end{enumerate}
\resultado \textbf{control predictivo lineal o cuadrático alimentado por predicciones corregidas, con salvaguardas por reglas independientes y monitorización del error de predicción y de los KPIs.}
\comprobacion el diseño respeta la separación entre restricciones duras (salvaguardas) y objetivo, y cada bloque tiene un indicador que permite detectar su fallo.
\end{solucion}

% ---------------------------------------------------------------------
\section{Tema 2}

\begin{solucion}{eje:2-holgura}
\datos óptimo $(100,100)$; holguras: molino 0, granuladora 0, demanda $160-100=60$.
\planteamiento \cref{cor:holgura}.
\pasos
\begin{enumerate}[label=\arabic*.]
  \item $y_3(160-100)=0$ con holgura $60>0$ obliga a $y_3^{*}=0$: un recurso que sobra no tiene valor marginal.
  \item $E=100>0$ obliga a que la restricción dual de $E$ sea activa: $0{,}5y_1+0{,}6y_2+0\cdot y_3=40$; igual con $P=100>0$: $0{,}5y_1+0{,}3y_2+y_3=30$.
  \item Precio máximo por una hora de granuladora: su precio sombra, $33{,}33$~€, válido mientras la base no cambie.
\end{enumerate}
\resultado \textbf{$y_3^{*}=0$ porque la demanda no está saturada; las duales son igualdades porque se fabrican ambos productos; se pagarían como mucho 33,33~€/h.}
\comprobacion el cuaderno resuelve con 91~h de granuladora y obtiene $7033{,}33$~€, exactamente $33{,}33$~€ más.
\end{solucion}

\begin{solucion}{eje:2-dual}
\datos primal $\max 40E+30P$ s.a. $0{,}5E+0{,}5P\le 100$, $0{,}6E+0{,}3P\le 90$, $P\le 160$, $E,P\ge 0$.
\planteamiento \cref{def:dual}: una variable por restricción, una restricción por variable.
\pasos
\begin{enumerate}[label=\arabic*.]
  \item Dual: $\min 100y_1+90y_2+160y_3$ s.a. $0{,}5y_1+0{,}6y_2\ge 40$, $0{,}5y_1+0{,}3y_2+y_3\ge 30$, $y\ge 0$.
  \item Con $y^{*}=(40;\ 100/3;\ 0)$: $20+20=40\ge 40$ y $20+10+0=30\ge 30$. Factible.
  \item $b\T y^{*}=4000+3000+0=7000$.
\end{enumerate}
\resultado \textbf{$y^{*}$ es dual factible y $b\T y^{*}=c\T x^{*}=7000$}: ambos son óptimos (\cref{teo:dual-debil}).
\comprobacion la holgura complementaria se cumple: las dos restricciones duales son activas y $E,P>0$; $y_3=0$ y la demanda tiene holgura.
\end{solucion}

\begin{solucion}{eje:2-outdoors}
\datos precios sombra $1{,}1667$ (curvado), 0 (soldadura), $0{,}80$ (tubo); aumento permitido del tubo: 555,56~kg; margen óptimo $2766{,}67$~\$.
\planteamiento \cref{prop:sombra}: dentro del rango, el valor óptimo cambia en precio sombra por unidad; un producto nuevo compensa si su margen supera el valor de los recursos que consume.
\pasos
\begin{enumerate}[label=\arabic*.]
  \item (a) 550~kg está dentro del rango (550 < 555,56). Aumento del margen bruto: $550\cdot 0{,}80=440$~\$. Coste: $550\cdot 0{,}70=385$~\$. Ganancia neta: $55$~\$. Conviene.
  \item (b) Valor de los recursos del toldo: $1{,}8\cdot 1{,}1667+0{,}5\cdot 0+1{,}3\cdot 0{,}80=2{,}10+1{,}04=3{,}14$~\$.
\end{enumerate}
\resultado \textbf{(a) sí: el margen sube 440~\$ y la ganancia neta es de 55~\$; (b) el toldo debe aportar más de 3,14~\$ por unidad.}
\comprobacion resolviendo con 2550~kg de tubo se obtiene $3206{,}67=2766{,}67+440$~\$ (480 sillas y 353,33 mesas).
\end{solucion}

\begin{solucion}{eje:2-bb}
\datos $\max x_1+5x_2$ s.a. $-4x_1+3x_2\le 6$, $3x_1+2x_2\le 18$, $x\in\Z^{2}_{\ge 0}$.
\planteamiento ramificación y acotación en profundidad, ramificando en una variable fraccionaria.
\pasos
\begin{enumerate}[label=\arabic*.]
  \item Raíz: el óptimo de la relajación está en la intersección de las dos rectas: de $-4x_1+3x_2=6$ y $3x_1+2x_2=18$ sale $x=(42/17,\ 90/17)=(2{,}47;\ 5{,}29)$, $z=28{,}94$. Se ramifica en $x_2$.
  \item $x_2\ge 6$: $-4x_1+18\le 6$ exige $x_1\ge 3$ y $3x_1+12\le 18$ exige $x_1\le 2$: infactible.
  \item $x_2\le 5$: óptimo en $x_2=5$, $x_1=8/3$ (sobre $3x_1+2x_2=18$), $z=27{,}67$. Se ramifica en $x_1$.
  \item $x_2\le 5$, $x_1\ge 3$: $x=(3;\ 4{,}5)$, $z=25{,}5$. Se ramifica en $x_2$: con $x_2\ge 5$ es infactible; con $x_2\le 4$, $x=(10/3;\ 4)$, $z=23{,}33$; se ramifica en $x_1$: $x_1\ge 4$ da $(4,3)$, $z=19$, primer incumbente; $x_1\le 3$ da $(3,4)$, $z=23$, nuevo incumbente.
  \item $x_2\le 5$, $x_1\le 2$: $x=(2;\ 4{,}67)$, $z=25{,}33$. Con $x_2\ge 5$ es infactible; con $x_2\le 4$, $(2,4)$, $z=22<23$: podado por cota.
\end{enumerate}
\resultado \textbf{óptimo $x^{*}=(3,4)$, $z^{*}=23$}, frente a $28{,}94$ de la relajación (11 nodos en profundidad).
\comprobacion $(3,4)$ cumple $-12+12=0\le 6$ y $9+8=17\le 18$; la enumeración de los enteros factibles confirma que ninguno supera 23.
\end{solucion}

\begin{solucion}{eje:2-mochila}
\datos valores $v=(33,45,25,17,39,23)$, pesos $w=(10,14,8,6,12,8)$, presupuesto 28.
\planteamiento la relajación de una mochila se resuelve llenando por orden de $v_j/w_j$ (Dantzig).
\pasos
\begin{enumerate}[label=\arabic*.]
  \item (a) $\max v\T x$ s.a. $w\T x\le 28$, $x\in\{0,1\}^{6}$.
  \item (b) Ratios: $3{,}30$; $3{,}21$; $3{,}13$; $2{,}83$; $3{,}25$; $2{,}88$. Orden: 1, 5, 2, 3, 6, 4. Se toman $x_1=1$ (capital 10) y $x_5=1$ (22); quedan 6 para la 2, que pide 14: $x_2=6/14=0{,}4286$. Valor: $33+39+45\cdot 0{,}4286=91{,}29$.
  \item (c) Como los valores son enteros, la cota útil es $\lfloor 91{,}29\rfloor=91$; gap $=(91{,}29-89)/89=2{,}6\,\%$, o $(91-89)/89=2{,}2\,\%$ con la cota entera.
  \item (d) Rama $x_2=1$: quedan 14 de capital; 1 entra entero (4 de capital restante), 5 entra en $4/12$: $45+33+13=91$. Rama $x_2=0$: 1 y 5 (22), quedan 6 para la 3: $x_3=0{,}75$, $33+39+18{,}75=90{,}75$.
  \item Ninguna rama puede podarse: $\lfloor 91\rfloor=91>89$ y $\lfloor 90{,}75\rfloor=90>89$.
\end{enumerate}
\resultado \textbf{relajación $91{,}29$ (cota 91); hijos $91$ y $90{,}75$; ninguno se poda.} El óptimo, 89, requiere 17 nodos en profundidad.
\comprobacion $\{1,4,5\}$ usa $10+6+12=28$ y vale $33+17+39=89$; la enumeración de las 64 carteras no encuentra ninguna mejor.
\end{solucion}

\begin{solucion}{eje:2-logica}
\datos $x_1,x_2\in\{0,1\}$.
\planteamiento cotas superiores e inferiores que fuercen el valor lógico.
\pasos
\begin{enumerate}[label=\arabic*.]
  \item Conjunción: $w\le x_1$, $w\le x_2$, $w\ge x_1+x_2-1$, $w\ge 0$.
  \item Disyunción: $y\ge x_1$, $y\ge x_2$, $y\le x_1+x_2$, $y\le 1$.
  \item O exclusivo con $a=x_1\wedge x_2$ (restricciones del paso~1): $z=x_1+x_2-2a$.
  \item Comprobación de $(x_1,x_2)$: $(0,0)$: $w=0$, $y=0$, $z=0$; $(1,0)$ y $(0,1)$: $w=0$, $y=1$, $z=1$; $(1,1)$: $w\ge 1$ luego $w=1$, $y=1$, $a=1$, $z=0$.
\end{enumerate}
\resultado \textbf{$w$: $w\le x_1$, $w\le x_2$, $w\ge x_1+x_2-1$; $y$: $y\ge x_1$, $y\ge x_2$, $y\le x_1+x_2$; $z=x_1+x_2-2a$ con $a=x_1\wedge x_2$.}
\comprobacion las cuatro combinaciones reproducen las tablas de verdad, y en cada caso las restricciones fijan un único valor entero.
\end{solucion}

\begin{solucion}{eje:2-bigm}
\datos $C=280$; demanda pendiente de Lager desde la semana~3: $160+190=350$; de IPA en la semana~4: 100.
\planteamiento $M_{pt}=\min(C,\sum_{s\ge t}d_{ps})$ (\cref{prop:bigm}); en la relajación, el menor $y$ compatible con $x\le My$ es $x/M$.
\pasos
\begin{enumerate}[label=\arabic*.]
  \item (a) Lager, semana~3: $\min(280,350)=280$. IPA, semana~4: $\min(280,100)=100$.
  \item (b) Con $x=170$ y $M=280$: $y\ge 0{,}607$, y se «paga» $0{,}607\cdot 500=303{,}6$~€ de preparación. Con $M=10^{4}$: $y\ge 0{,}017$ y se pagan $8{,}50$~€.
  \item (c) $y_{\mathrm{IPA},t}+y_{\mathrm{Tost},t}\le 1$ para cada $t$.
\end{enumerate}
\resultado \textbf{(a) 280 y 100~hL; (b) $y\ge 0{,}607$ (303,6~€) frente a $y\ge 0{,}017$ (8,5~€); (c) $y_{\mathrm{IPA},t}+y_{\mathrm{Tost},t}\le 1$.}
\comprobacion con la $M$ grande la relajación casi ignora las preparaciones y su cota empeora, como predice la \cref{prop:bigm}~(ii).
\end{solucion}

\begin{solucion}{eje:2-concava}
\datos PL $\max\{c\T x: Ax\le b,\ x\ge 0\}$ con poliedro dual $D=\{y\ge 0: A\T y\ge c\}$ no vacío.
\planteamiento por dualidad fuerte, $z^{*}(b)=\min_{y\in D}b\T y$.
\pasos
\begin{enumerate}[label=\arabic*.]
  \item Si el primal es factible para $b$ y el dual es factible, el primal está acotado (\cref{teo:dual-debil}) y tiene óptimo; por la \cref{teo:dual-fuerte}, $z^{*}(b)=\min_{y\in D}b\T y$.
  \item Si $D$ tiene vértices $y^{1},\dots,y^{K}$ y el mínimo se alcanza en un vértice (\cref{teo:vertice} aplicado al dual), $z^{*}(b)=\min_kb\T y^{k}$.
  \item Un mínimo de funciones lineales de $b$ es cóncavo y lineal a trozos.
  \item Para un recurso $i$, la función $\beta\mapsto z^{*}(b+\beta e_i)$ es cóncava en $\beta$, así que su pendiente (el precio sombra $y_i^{*}$) es no creciente en $\beta$.
\end{enumerate}
\resultado \textbf{$z^{*}(b)=\min_ky^{k\mathsf{T}}b$ es cóncava y lineal a trozos; el valor marginal de cada recurso no crece al aumentar su disponibilidad.}
\comprobacion en la cervecera, la capacidad de la semana~4 vale 4~€/hL hasta 290~hL y 0 después: pendiente no creciente.
\end{solucion}

% ---------------------------------------------------------------------
\section{Tema 3}

\begin{solucion}{eje:3-contar}
\datos $n=10$ y $n=20$; $10^{9}$ circuitos por segundo.
\planteamiento \cref{prop:circuitos}: $(n-1)!/2$.
\pasos
\begin{enumerate}[label=\arabic*.]
  \item $n=10$: $9!/2=362\,880/2=181\,440$ circuitos; tiempo $1{,}8\cdot 10^{-4}$~s.
  \item $n=20$: $19!/2=1{,}216\cdot 10^{17}/2=6{,}08\cdot 10^{16}$; tiempo $6{,}08\cdot 10^{7}$~s, unos 1{,}9 años.
\end{enumerate}
\resultado \textbf{181\,440 circuitos (0,18~ms) y $6{,}08\cdot 10^{16}$ circuitos (casi 2 años).}
\comprobacion la tabla de 15.053, con $10^{21}$ pasos por segundo, muestra el mismo crecimiento: $n!$ pasa de milisegundos con 20 ciudades a miles de años con 30.
\end{solucion}

\begin{solucion}{eje:3-ruta}
\datos $t_{01}=16$, $t_{12}=10$, $t_{24}=22$, $t_{43}=14$, $t_{30}=14$; $s=10$; ventanas del \cref{ej:vrptw}.
\planteamiento llegada $=\max(a_j,\ \text{salida}+t_{ij})$, con salida $=$ llegada $+\,s$.
\pasos
\begin{enumerate}[label=\arabic*.]
  \item Bar 1: llega a los 16, espera hasta 60 (44~min), sale a los 70.
  \item Bar 2: $70+10=80\le 90$; sale a los 90.
  \item Bar 4: $90+22=112\in[30,120]$; sale a los 122.
  \item Bar 3: $122+14=136\in[120,180]$; sale a los 146.
  \item Almacén: $146+14=160\le 300$. Carga: $4+6+3+5=18\le 22$. Viaje: $16+10+22+14+14=76$~min.
  \item Ruta $0\to 3\to 1$: llega al bar~3 a los 14, espera hasta 120, sale a los 130 y llega al bar~1 a los $130+22=152>120$: infactible.
\end{enumerate}
\resultado \textbf{la primera ruta es factible (llegadas 60, 80, 112, 136 y regreso a los 160; 76~min de viaje); la segunda no.}
\comprobacion coincide con \texttt{evalua\_ruta} del cuaderno.
\end{solucion}

\begin{solucion}{eje:3-2opt}
\datos $A=(0,0)$, $B=(4,0)$, $C=(4,3)$, $D=(0,3)$.
\planteamiento \cref{def:2opt}; las aristas $(A,C)$ y $(B,D)$ son las diagonales y se cruzan.
\pasos
\begin{enumerate}[label=\arabic*.]
  \item $A\to C\to B\to D\to A$: $5+3+5+3=16$.
  \item El 2-opt elimina $(A,C)$ y $(B,D)$, añade $(A,B)$ y $(C,D)$ e invierte el tramo: $A\to B\to C\to D\to A$, longitud $4+3+4+3=14$.
  \item (b) Si $(a,b)$ y $(c,d)$ se cortan en $P$, por la desigualdad triangular $d(a,c)\le d(a,P)+d(P,c)$ y $d(b,d)\le d(b,P)+d(P,d)$. Sumando, $d(a,c)+d(b,d)\le d(a,b)+d(c,d)$, con desigualdad estricta salvo si los puntos están alineados.
\end{enumerate}
\resultado \textbf{de 16 a 14 (ahorro de 2); deshacer un cruce nunca alarga el circuito.}
\comprobacion $\Delta=d(A,B)+d(C,D)-d(A,C)-d(B,D)=4+4-5-5=-2<0$.
\end{solucion}

\begin{solucion}{eje:3-cuñas}
\datos duraciones 12, 18, 22, 35, 40, 59~s; bloques de 60~s.
\planteamiento empaquetado en contenedores: $y_k=1$ si se usa el bloque $k$; $x_{ik}=1$ si el anuncio $i$ va en el bloque $k$.
\pasos
\begin{enumerate}[label=\arabic*.]
  \item PLEM: $\min\sum_ky_k$ s.a. $\sum_kx_{ik}=1$ para cada anuncio; $\sum_i\ell_ix_{ik}\le 60y_k$ para cada bloque; binarias, con $k=1,\dots,6$.
  \item Cota: $\sum\ell_i=186$; $\lceil 186/60\rceil=4$.
  \item Solución con 4 bloques: $\{59\}$, $\{40,18\}$, $\{35,22\}$, $\{12\}$ (o $\{40,12\}$, $\{35,22\}$, $\{59\}$, $\{18\}$).
\end{enumerate}
\resultado \textbf{4 bloques, óptimo por coincidir con la cota.}
\comprobacion $40+18=58$, $35+22=57$; ningún par cabe con el anuncio de 59~s.
\end{solucion}

\begin{solucion}{eje:3-johnson}
\datos $J1=(3,6)$, $J2=(5,2)$, $J3=(1,2)$, $J4=(6,6)$, $J5=(7,5)$.
\planteamiento regla de Johnson y recurrencia $C_1\leftarrow C_1+p_1$, $C_2\leftarrow\max(C_2,C_1)+p_2$.
\pasos
\begin{enumerate}[label=\arabic*.]
  \item $p_1\le p_2$: J1, J3, J4; por $p_1$ creciente: J3, J1, J4. $p_1>p_2$: J2, J5; por $p_2$ decreciente: J5, J2. Secuencia: J3, J1, J4, J5, J2.
  \item J3: $C_1=1$, $C_2=3$. J1: $C_1=4$, $C_2=\max(3,4)+6=10$. J4: $C_1=10$, $C_2=16$. J5: $C_1=17$, $C_2=\max(16,17)+5=22$. J2: $C_1=22$, $C_2=\max(22,22)+2=24$.
  \item Cota de la \cref{prop:cota-ms}: cargas $22$ y $21$, piezas de a lo sumo 13~h: $C_{\max}\ge 22$.
\end{enumerate}
\resultado \textbf{secuencia J3, J1, J4, J5, J2 con makespan 24~h}; cota 22~h.
\comprobacion una cota más fina es la carga de la sierra más el menor $p_2$ de la última pieza, $22+2=24$, de modo que 24 es óptimo; la enumeración de las 120 secuencias lo confirma.
\end{solucion}

\begin{solucion}{eje:3-bigm}
\datos $b_1=120$, $b_8=210$, $b_3=180$; $a_2=0$, $a_5=150$, $a_1=60$; $t_{12}=10$, $t_{85}=24$, $t_{31}=22$; $s=10$.
\planteamiento $M_{ij}=\max(0,\ b_i+s+t_{ij}-a_j)$ (\cref{prop:bigm-tw}); un arco es imposible si $a_i+s+t_{ij}>b_j$.
\pasos
\begin{enumerate}[label=\arabic*.]
  \item $M_{12}=120+10+10-0=140$.
  \item $M_{85}=210+10+24-150=94$.
  \item $M_{31}=180+10+22-60=152$.
  \item El bar~3 no puede servirse antes de los 120~min; la llegada más temprana al bar~1 por ese arco es $120+10+22=152>120=b_1$. Ninguna ruta puede usar $3\to 1$.
  \item Preproceso: eliminar todo arco con $a_i+s+t_{ij}>b_j$; en este caso hay 19 de los 56 arcos entre clientes.
\end{enumerate}
\resultado \textbf{$M_{12}=140$, $M_{85}=94$, $M_{31}=152$; el arco $3\to 1$ es infactible y se elimina, como otros 18.}
\comprobacion la ruta óptima no usa ningún arco eliminado: sus arcos entre clientes son $1\to 2$, $2\to 4$, $4\to 3$, $6\to 7$, $7\to 8$ y $8\to 5$.
\end{solucion}

\begin{solucion}{eje:3-cota}
\datos caso del \cref{ej:jobshop}; cota 11~h.
\planteamiento (a) como en la \cref{prop:cota-ms}; (b) construcción; (c) análisis de casos suponiendo $C_{\max}=11$.
\pasos
\begin{enumerate}[label=\arabic*.]
  \item (a) Véase la demostración de la \cref{prop:cota-ms}.
  \item (b) Se alcanza: J1 $=(M_1,2),(M_2,1)$, J2 $=(M_2,2),(M_1,1)$: se programa J1 en $M_1$ [0,2] y J2 en $M_2$ [0,2], después J1 en $M_2$ [2,3] y J2 en $M_1$ [2,3]; $C_{\max}=3=$ cota. No se alcanza: J1 $=(M_1,2),(M_2,2)$ y J2 $=(M_1,2),(M_2,2)$: la cota es 4, pero $M_2$ no puede empezar antes de las 2 y debe trabajar 4~h, así que $C_{\max}=6$.
  \item (c) Con $C_{\max}=11$ la fresadora (11~h de carga) trabaja sin pausa de 0 a 11. En $t=0$ solo puede empezar P2 (las demás operaciones de fresadora van precedidas de otras): P2 ocupa [0,4].
  \item En $t=4$ debe empezar otra: P3 y P4 necesitan antes $3+3=6$ y $2+3=5$~h, luego solo P1, que exige su torno (3~h) terminado: torno con P1 en [0,3] o [1,4]; fresadora con P1 en [4,6].
  \item En $t=6$: P4 necesitaría torno y taladro terminados, pero el torno está ocupado por P1 hasta las 3 como pronto, y $3+2+3=8>6$. Luego P3, que necesita su taladro [0,3] y su torno de 3~h terminado a las 6: torno con P3 en [3,6], lo que obliga a P1 en [0,3]. Fresadora con P3 en [6,8].
  \item En $t=8$, P4: su torno solo puede ir en [6,8] (el torno está ocupado hasta las 6) y su taladro en [8,11], de modo que su fresado empezaría a las 11. Contradicción.
\end{enumerate}
\resultado \textbf{la cota de 11~h es inalcanzable y $C_{\max}^{*}\ge 12$}; el programa del \cref{ej:jobshop} alcanza 12, luego es óptimo.
\comprobacion CP-SAT y la enumeración de los 13\,824 órdenes del cuaderno dan 12~h.
\end{solucion}

% ---------------------------------------------------------------------
\section{Tema 4}

\begin{solucion}{eje:4-dcp}
\datos cinco problemas.
\planteamiento \cref{def:convexo}: objetivo convexo, desigualdades $g\le 0$ con $g$ convexa, igualdades afines.
\pasos
\begin{enumerate}[label=\arabic*.]
  \item (a) Convexo: cuadrática definida positiva y restricción lineal.
  \item (b) Convexo: $-\log$ es convexa y la igualdad es afín.
  \item (c) No convexo: la igualdad $x_1^{2}+x_2^{2}=1$ no es afín (el conjunto es una circunferencia).
  \item (d) Convexo aunque no lo parezca: $\max$ es convexa, y $\{x\ge 0: x_1x_2\ge 1\}$ es convexo (es $\{x>0: \log x_1+\log x_2\ge 0\}$, superconjunto de nivel de una función cóncava). Para DCP hay que escribirla como \texttt{geo\_mean(x) >= 1}.
  \item (e) Convexo: norma de una función afín y restricción de norma.
\end{enumerate}
\resultado \textbf{convexos (a), (b), (d) y (e); no convexo (c).}
\comprobacion (c) tiene dos mínimos en $(0,\pm 1)$, ambos globales, pero el conjunto no es convexo: el punto medio $(0,0)$ no es factible.
\end{solucion}

\begin{solucion}{eje:4-kkt}
\datos $\min x_1^{2}+x_2^{2}$ s.a. $5-x_1-2x_2\le 0$.
\planteamiento \cref{teo:kkt} con $g(x)=5-x_1-2x_2$.
\pasos
\begin{enumerate}[label=\arabic*.]
  \item Estacionariedad: $(2x_1,2x_2)+\lambda(-1,-2)=0$, de donde $x=\tfrac{\lambda}{2}(1,2)$.
  \item Si $\lambda=0$, $x=0$, infactible. Luego la restricción es activa: $\tfrac{\lambda}{2}(1+4)=5$, $\lambda=2$, $x=(1,2)$, $p^{*}=5$.
  \item Sensibilidad: con lado derecho $u$, $x=\tfrac{u}{5}(1,2)$ y $p^{*}(u)=u^{2}/5$; $\mathrm{d}p^{*}/\mathrm{d}u=2u/5=2=\lambda$ en $u=5$.
  \item Estimación: $\Delta p^{*}\approx 2\cdot 0{,}1=0{,}2$. Exacto: $5{,}1^{2}/5-5=5{,}202-5=0{,}202$.
\end{enumerate}
\resultado \textbf{$x^{*}=(1,2)$, $\lambda^{*}=2$; el valor óptimo sube unas 0,2 unidades (exactamente 0,202).}
\comprobacion $x^{*}$ es la proyección del origen sobre la recta $x_1+2x_2=5$: distancia al cuadrado $25/5=5$.
\end{solucion}

\begin{solucion}{eje:4-despacho2}
\datos $C_1=20P_1+0{,}01P_1^{2}$, $C_2=18P_2+0{,}02P_2^{2}$, $D=500$~MW.
\planteamiento igual coste incremental: $20+0{,}02P_1=18+0{,}04P_2=\lambda$ y $P_1+P_2=500$.
\pasos
\begin{enumerate}[label=\arabic*.]
  \item $P_1=50(\lambda-20)$, $P_2=25(\lambda-18)$.
  \item $50\lambda-1000+25\lambda-450=500$, luego $75\lambda=1950$ y $\lambda=26$~€/MWh.
  \item $P_1=300$~MW, $P_2=200$~MW.
  \item Coste: $C_1=6000+900=6900$; $C_2=3600+800=4400$; total $11\,300$~€/h.
\end{enumerate}
\resultado $\bm{P=(300,200)}$~MW, $\bm{\lambda=26}$~€/MWh, coste \textbf{11\,300~€/h}.
\comprobacion repartos vecinos son más caros: con $(301,199)$, $C_1=6926{,}01$ y $C_2=4374{,}02$, total $11\,300{,}03$.
\end{solucion}

\begin{solucion}{eje:4-agua}
\datos $\max\sum_i\log(\alpha_i+x_i)$, $\alpha=(0{,}5;\ 1;\ 1{,}5)$, $\sum x_i=1$, $x\ge 0$.
\planteamiento KKT del problema equivalente $\min-\sum\log(\alpha_i+x_i)$ con multiplicadores $\nu$ (igualdad) y $\lambda_i\ge 0$ ($-x_i\le 0$).
\pasos
\begin{enumerate}[label=\arabic*.]
  \item Estacionariedad: $-1/(\alpha_i+x_i)-\lambda_i+\nu=0$; holgura: $\lambda_ix_i=0$.
  \item Si $x_i>0$: $\lambda_i=0$ y $x_i=1/\nu-\alpha_i$. Si $x_i=0$: $\lambda_i=\nu-1/\alpha_i\ge 0$, es decir, $1/\nu\le\alpha_i$. En ambos casos $x_i=\max(0,\ 1/\nu-\alpha_i)$.
  \item Probando los tres canales activos: $3/\nu-3=1$, $1/\nu=4/3$, $x_3=4/3-1{,}5<0$: no vale.
  \item Con los canales 1 y 2: $2/\nu-1{,}5=1$, $1/\nu=1{,}25$; $x=(0{,}75;\ 0{,}25;\ 0)$, y $1{,}25\le 1{,}5$ confirma $x_3=0$.
\end{enumerate}
\resultado $\bm{\nu=0{,}8}$, $\bm{x^{*}=(0{,}75;\ 0{,}25;\ 0)}$.
\comprobacion $0{,}75+0{,}25=1$; el «nivel del agua» $\alpha_i+x_i$ es $1{,}25$ en los canales activos y el canal~3 ($\alpha_3=1{,}5$) queda por encima.
\end{solucion}

\begin{solucion}{eje:4-sensib}
\datos $\lambda=28{,}4$; $P_2=350$ en su máximo; incrementales $18+0{,}04P_1$, $20+0{,}02P_2$, $26+0{,}01P_3$.
\planteamiento (a) \cref{prop:sensibilidad} y nuevo cálculo; (b) mientras se mantenga la estructura de límites activos.
\pasos
\begin{enumerate}[label=\arabic*.]
  \item (a) Estimación: $20\cdot 28{,}4=568$~€/h.
  \item Exacto con $D=830$: $25(\lambda-18)+100(\lambda-26)=480$, $\lambda=28{,}24$, $P=(256,350,224)$, coste $21\,418{,}6$~€/h; ahorro $21\,985-21\,418{,}6=566{,}4$~€/h.
  \item (b) Al subir $D$, el primero en tocar un límite es el grupo~3 con $26+0{,}01P_3=\lambda$ y $P_3=300$, es decir, $\lambda=29$, $P_1=275$: $D=275+350+300=925$.
  \item Al bajar $D$, el grupo~2 deja su máximo cuando $\lambda=20+0{,}02\cdot 350=27$: $P_1=225$, $P_3=100$, $D=675$.
\end{enumerate}
\resultado \textbf{(a) ahorro estimado 568~€/h, exacto 566,4~€/h; (b) para $675\le D\le 925$~MW solo está activo el máximo del grupo~2.}
\comprobacion en ambos extremos los mínimos técnicos no se tocan ($P_1\ge 225>100$, $P_3\ge 100>50$).
\end{solucion}

\begin{solucion}{eje:4-dual}
\datos $C_i=a_i+b_iP_i+c_iP_i^{2}$ con los datos del \cref{ej:despacho}, sin límites; $D=850$.
\planteamiento minimizar la lagrangiana en cada $P_i$ por separado.
\pasos
\begin{enumerate}[label=\arabic*.]
  \item $\min_{P_i}\bigl[a_i+(b_i-\lambda)P_i+c_iP_i^{2}\bigr]$ se alcanza en $P_i=(\lambda-b_i)/(2c_i)$ y vale $a_i-(\lambda-b_i)^{2}/(4c_i)$.
  \item $q(\lambda)=\sum_ia_i-\sum_i(\lambda-b_i)^{2}/(4c_i)+\lambda D$.
  \item $q'(\lambda)=-\sum_i(\lambda-b_i)/(2c_i)+D=0$ es la ecuación de balance: $\lambda=28$.
  \item $q(28)=1200-\bigl(10^{2}/0{,}08+8^{2}/0{,}04+2^{2}/0{,}02\bigr)+28\cdot 850=1200-(1250+1600+200)+23\,800=21\,950$.
  \item Primal en $P=(250,400,200)$: $6250+10\,000+5700=21\,950$.
\end{enumerate}
\resultado \textbf{$q(\lambda)=1200-\sum_i(\lambda-b_i)^{2}/(4c_i)+850\lambda$, máximo en $\lambda=28$ con valor 21\,950~€/h, igual al coste primal}: dualidad fuerte.
\comprobacion $q''(\lambda)=-\sum_i1/(2c_i)=-175<0$: el punto crítico es un máximo.
\end{solucion}

\begin{solucion}{eje:4-primer-orden}
\datos $f$ convexa y diferenciable, $\mathcal{X}$ convexo.
\planteamiento desigualdad de primer orden de la convexidad, $f(x)\ge f(x^{*})+\nabla f(x^{*})\T(x-x^{*})$ (módulo~4, tema~5).
\pasos
\begin{enumerate}[label=\arabic*.]
  \item Suficiencia: si $\nabla f(x^{*})\T(x-x^{*})\ge 0$ para todo $x\in\mathcal{X}$, la desigualdad de primer orden da $f(x)\ge f(x^{*})$.
  \item Necesidad: si existiera $x\in\mathcal{X}$ con $\nabla f(x^{*})\T(x-x^{*})<0$, la derivada direccional de $f$ en $x^{*}$ hacia $x$ sería negativa y, para $t>0$ pequeño, $x^{*}+t(x-x^{*})\in\mathcal{X}$ (convexidad) tendría $f$ menor que $f(x^{*})$.
  \item Geometría: $-\nabla f(x^{*})$ forma un ángulo de al menos $90^{\circ}$ con toda dirección factible; el gradiente negativo «apunta hacia fuera» de $\mathcal{X}$.
  \item Si $x^{*}$ es interior, las direcciones $\pm e_k$ son factibles: $\pm\partial_kf(x^{*})\ge 0$ implica $\nabla f(x^{*})=0$.
\end{enumerate}
\resultado \textbf{$x^{*}$ es óptimo si y solo si $\nabla f(x^{*})\T(x-x^{*})\ge 0$ para todo $x\in\mathcal{X}$; en el interior, $\nabla f(x^{*})=0$.}
\comprobacion en el \cref{eje:4-kkt}, $\nabla f(1,2)=(2,4)$ y para $x$ factible $(2,4)\T(x-(1,2))=2(x_1+2x_2-5)\ge 0$.
\end{solucion}

% ---------------------------------------------------------------------
\section{Tema 5}

\begin{solucion}{eje:5-aceptacion}
\datos $\Delta=50$; $T=100$ y $T=25$; $T_{k+1}=0{,}99T_k$, $T_0=500$.
\planteamiento probabilidad $e^{-\Delta/T}$; $500\cdot 0{,}99^{k}<1$.
\pasos
\begin{enumerate}[label=\arabic*.]
  \item $T=100$: $e^{-0{,}5}=0{,}6065$.
  \item $T=25$: $e^{-2}=0{,}1353$.
  \item $k>\ln 500/(-\ln 0{,}99)=6{,}2146/0{,}010050=618{,}3$, luego $k=619$.
\end{enumerate}
\resultado \textbf{60,65\,\% y 13,53\,\%; hacen falta 619 iteraciones.}
\comprobacion $500\cdot 0{,}99^{619}=0{,}993<1$ y $500\cdot 0{,}99^{618}=1{,}004>1$.
\end{solucion}

\begin{solucion}{eje:5-frentes}
\datos $a=(1,9)$, $b=(2,7)$, $c=(3,8)$, $d=(4,4)$, $e=(5,5)$, $f=(7,2)$.
\planteamiento \cref{def:pareto}: se buscan los no dominados, se retiran y se repite.
\pasos
\begin{enumerate}[label=\arabic*.]
  \item $c$ está dominado por $b$ ($2\le 3$, $7\le 8$); $e$ por $d$ ($4\le 5$, $4\le 5$).
  \item $a$, $b$, $d$ y $f$ no están dominados: cada uno es el mejor en $f_1$ entre los que tienen su $f_2$ o menos.
  \item Sin el primer frente, $c$ y $e$ no se dominan entre sí ($3<5$ pero $8>5$).
\end{enumerate}
\resultado \textbf{$\mathcal{F}_1=\{a,b,d,f\}$ (rango 1) y $\mathcal{F}_2=\{c,e\}$ (rango 2).}
\comprobacion la función \texttt{ordenacion\_no\_dominada} del cuaderno devuelve los mismos rangos.
\end{solucion}

\begin{solucion}{eje:5-crowding}
\datos $\mathcal{F}_1$: $a=(1,9)$, $b=(2,7)$, $d=(4,4)$, $f=(7,2)$; rangos $f_1^{\max}-f_1^{\min}=6$, $f_2^{\max}-f_2^{\min}=7$.
\planteamiento distancia de aglomeración de NSGA-II.
\pasos
\begin{enumerate}[label=\arabic*.]
  \item Orden por $f_1$: $a,b,d,f$; extremos $a$ y $f$: distancia infinita.
  \item $b$: $(4-1)/6+(9-4)/7=0{,}5+0{,}7143=1{,}214$ (en $f_2$ sus vecinos son $d$ y $a$).
  \item $d$: $(7-2)/6+(7-2)/7=0{,}8333+0{,}7143=1{,}548$.
\end{enumerate}
\resultado \textbf{$a,f$: $\infty$; $b$: 1,214; $d$: 1,548. Se descartaría primero $b$}, el más rodeado.
\comprobacion la función \texttt{crowding} del cuaderno da los mismos valores con estos cuatro puntos.
\end{solucion}

\begin{solucion}{eje:5-apoyado}
\datos $A=(3400;\ 13{,}5)$, $B=(3700;\ 10{,}5)$, $C=(3900;\ 7{,}833)$, $D=(4400;\ 6{,}833)$.
\planteamiento para cada $w$, comparar $B$ con $A$ y con $C$.
\pasos
\begin{enumerate}[label=\arabic*.]
  \item $A$ es mejor que $B$ si $w\cdot 3400+(1-w)\,13{,}5<w\cdot 3700+(1-w)\,10{,}5$, es decir, $3(1-w)<300w$: $w>3/303=0{,}00990$.
  \item $C$ es mejor que $B$ si $200w<2{,}667(1-w)$: $w<2{,}667/202{,}67=0{,}01316$.
  \item Como $0{,}00990<0{,}01316$, todo $w\in[0,1]$ cumple alguna de las dos condiciones.
\end{enumerate}
\resultado \textbf{para todo $w$, $A$ o $C$ tiene menor suma ponderada que $B$}: $B$ no está apoyado.
\comprobacion geométricamente, $B$ queda $0{,}4$ por encima del segmento $AC$ en $f_1=3700$.
\end{solucion}

\begin{solucion}{eje:5-ox}
\datos $p_1=(1,2,3,4,5,6)$, $p_2=(4,6,2,1,5,3)$, tramo de las posiciones 3 a~5.
\planteamiento copiar el tramo y rellenar con el orden de $p_2$.
\pasos
\begin{enumerate}[label=\arabic*.]
  \item Tramo copiado: $(\_,\_,3,4,5,\_)$.
  \item Genes de $p_2$ que faltan, en su orden: $6,2,1$ (se saltan 4, 5 y 3).
  \item Relleno de izquierda a derecha: posiciones 1, 2 y 6 reciben 6, 2 y 1.
\end{enumerate}
\resultado \textbf{hijo $(6,2,3,4,5,1)$.}
\comprobacion contiene cada número del 1 al 6 exactamente una vez.
\end{solucion}

\begin{solucion}{eje:5-pesos}
\datos $\lambda=6$, $\mu=3$.
\planteamiento $w_i'=\ln 3{,}5-\ln i$, normalizados.
\pasos
\begin{enumerate}[label=\arabic*.]
  \item $w'=(1{,}2528;\ 0{,}5596;\ 0{,}1542)$, suma $1{,}9665$.
  \item $w=(0{,}6370;\ 0{,}2846;\ 0{,}0784)$.
  \item $\sum w_i^{2}=0{,}4058+0{,}0810+0{,}0061=0{,}4929$; $\mu_{\mathrm{eff}}=2{,}029$.
\end{enumerate}
\resultado $\bm{w=(0{,}637;\ 0{,}285;\ 0{,}078)}$, $\bm{\mu_{\mathrm{eff}}\approx 2{,}03}$.
\comprobacion $1\le\mu_{\mathrm{eff}}\le\mu=3$, con igualdad superior solo para pesos iguales.
\end{solucion}

\begin{solucion}{eje:5-eps}
\datos problema $\min f_1$ s.a. $f_2\le\varepsilon$ con solución única $x^{*}$; frente del \cref{ej:planta}.
\planteamiento reducción al absurdo; elegir $\varepsilon=10{,}5$.
\pasos
\begin{enumerate}[label=\arabic*.]
  \item (a) Si $x$ dominara a $x^{*}$: $f_2(x)\le f_2(x^{*})\le\varepsilon$, así que $x$ es factible, y $f_1(x)\le f_1(x^{*})$. Por la unicidad, $x=x^{*}$, que no puede dominarse a sí mismo.
  \item (b) Con $\varepsilon=10{,}5$, las distribuciones con $f_2\le 10{,}5$ tienen $f_1\ge 3700$ (si alguna tuviera menos, dominaría a $(3700;\ 10{,}5)$ o tendría $f_2>10{,}5$), y el mínimo es $3700$.
\end{enumerate}
\resultado \textbf{(a) la solución única es Pareto-óptima; (b) $\varepsilon=10{,}5$ devuelve $(3700;\ 10{,}5)$.}
\comprobacion en el cuaderno, filtrando las 720 distribuciones con $f_2\le 10{,}5$, el menor $f_1$ es 3700.
\end{solucion}

\begin{solucion}{eje:5-metropolis}
\datos propuesta $\Prob(x\to x')$ no simétrica; aceptación de Metropolis--Hastings.
\planteamiento comprobar $\pi_T(x)P(x,x')=\pi_T(x')P(x',x)$.
\pasos
\begin{enumerate}[label=\arabic*.]
  \item Sea $r=\dfrac{\pi_T(x')\Prob(x'\to x)}{\pi_T(x)\Prob(x\to x')}$; la aceptación es $\alpha(x,x')=\min(1,r)$ y $\alpha(x',x)=\min(1,1/r)$.
  \item Si $r\le 1$: $\pi_T(x)\Prob(x\to x')\,r=\pi_T(x')\Prob(x'\to x)=\pi_T(x')\Prob(x'\to x)\cdot 1$.
  \item Si $r>1$: simétrico intercambiando papeles.
  \item Sin la corrección, el flujo de $x$ a $x'$ es proporcional a $\Prob(x\to x')$, y la cadena pasa más tiempo donde la propuesta lleva más masa, aunque $f$ no sea mejor allí.
\end{enumerate}
\resultado \textbf{con la corrección se cumple el equilibrio detallado con $\pi_T$; sin ella, la distribución estacionaria se sesga hacia las regiones favorecidas por la propuesta.}
\comprobacion con propuesta simétrica el cociente de propuestas vale~1 y se recupera la regla de la \cref{prop:boltzmann}.
\end{solucion}

% ---------------------------------------------------------------------
\section{Tema 6}

\begin{solucion}{eje:6-fugas}
\datos cuatro prácticas.
\planteamiento preguntar en cada caso qué información no estaría disponible al decidir o qué se usa dos veces.
\pasos
\begin{enumerate}[label=\arabic*.]
  \item (a) Residuos dentro de muestra de un modelo flexible: casi nulos, cuantiles demasiado estrechos y producción insuficiente. Hay que usar residuos de un periodo de calibración.
  \item (b) La temperatura real de la semana no se conoce el lunes: se debe usar la previsión disponible en ese momento.
  \item (c) Maldición del optimizador: el beneficio del plan elegido está sesgado al alza; hay que reevaluarlo con datos nuevos.
  \item (d) En series temporales, mezclar el futuro con el pasado en la validación permite que el modelo aprenda la tendencia futura; la validación debe respetar el orden temporal.
\end{enumerate}
\resultado \textbf{(a) calibración dentro de muestra; (b) variable no disponible al decidir; (c) sesgo de selección; (d) validación no temporal.}
\comprobacion el protocolo del \cref{ej:fresca} (entrenamiento, calibración y prueba consecutivos) evita las cuatro.
\end{solucion}

\begin{solucion}{eje:6-newsvendor}
\datos $D\sim\mathcal N(80,16^{2})$; $c_u=230$, $c_o=50$; después $c_o=70$.
\planteamiento \cref{teo:fractil}: $q^{*}=\mu+\sigma\Phi^{-1}(c_u/(c_u+c_o))$.
\pasos
\begin{enumerate}[label=\arabic*.]
  \item Fractil: $230/280=0{,}8214$; $\Phi^{-1}(0{,}8214)=0{,}9208$.
  \item $q^{*}=80+16\cdot 0{,}9208=94{,}73$~hL, 14,7~hL por encima de la media.
  \item Con saldo nulo, $c_o=70$: fractil $230/300=0{,}7667$, $\Phi^{-1}=0{,}7279$, $q^{*}=80+11{,}65=91{,}65$~hL.
\end{enumerate}
\resultado $\bm{q^{*}=94{,}7}$~\textbf{hL}; sin saldo, $\bm{91{,}6}$~\textbf{hL}.
\comprobacion al encarecerse el sobrante, $q^{*}$ baja, como debe; ambos superan la media porque el fractil es mayor que $0{,}5$.
\end{solucion}

\begin{solucion}{eje:6-spo}
\datos $c=(1{,}0;\ 1{,}2)$; $\hat c_A=(1{,}25;\ 1{,}2)$; $\hat c_B=(0{,}75;\ 1{,}2)$.
\planteamiento error cuadrático $\norm{\hat c-c}^{2}$; decisión $\argmin_k\hat c_k$; pérdida $c_{\text{elegida}}-\min_kc_k$.
\pasos
\begin{enumerate}[label=\arabic*.]
  \item Errores: A, $0{,}25^{2}=0{,}0625$; B, $(-0{,}25)^{2}=0{,}0625$.
  \item A elige la ruta~2 ($1{,}2<1{,}25$): pérdida $1{,}2-1{,}0=0{,}2$~h.
  \item B elige la ruta~1 ($0{,}75<1{,}2$): pérdida 0.
\end{enumerate}
\resultado \textbf{mismo error cuadrático (0,0625), pérdidas de decisión 0,2~h y 0~h.} Lo que importa es el orden que induce la predicción, no su exactitud media.
\comprobacion en el caso lineal, la decisión solo depende de la predicción a través del $\argmin$ (\cref{prop:lineal}); SPO mide precisamente esa pérdida.
\end{solucion}

\begin{solucion}{eje:6-vss}
\datos $D\in\{60,100,140\}$ con probabilidades $(0{,}3;\ 0{,}4;\ 0{,}3)$; $p=300$, $c=70$, $s=20$; $c_u=230$, $c_o=50$.
\planteamiento beneficio $\pi(q,D)=300\min(q,D)-70q+20(q-D)^{+}$; en maximización, $\mathrm{VSS}=\mathrm{RP}-\mathrm{EEV}$ y $\mathrm{EVPI}=\mathrm{WS}-\mathrm{RP}$.
\pasos
\begin{enumerate}[label=\arabic*.]
  \item (a) $\E D=100$, luego $\bar q=100$. Beneficios: $D=60$: $18\,000-7000+800=11\,800$; $D\ge 100$: $30\,000-7000=23\,000$. $\mathrm{EEV}=0{,}3\cdot 11\,800+0{,}7\cdot 23\,000=19\,640$~€.
  \item (b) Fractil $0{,}8214$; $F(100)=0{,}7<0{,}8214\le F(140)=1$: $q^{*}=140$. Beneficios: $D=60$: $18\,000-9800+1600=9800$; $D=100$: $30\,000-9800+800=21\,000$; $D=140$: $42\,000-9800=32\,200$. $\mathrm{RP}=2940+8400+9660=21\,000$~€.
  \item (c) $\mathrm{WS}=230\cdot\E D=230\cdot 100=23\,000$~€ (se produce exactamente $D$).
  \item $\mathrm{VSS}=21\,000-19\,640=1360$~€; $\mathrm{EVPI}=23\,000-21\,000=2000$~€.
\end{enumerate}
\resultado \textbf{EEV $=19\,640$, RP $=21\,000$ ($q=140$), WS $=23\,000$; VSS $=1360$~€ y EVPI $=2000$~€.}
\comprobacion se cumple $\mathrm{EEV}\le\mathrm{RP}\le\mathrm{WS}$ (\cref{prop:vss} en versión de maximización).
\end{solucion}

\begin{solucion}{eje:6-robusta}
\datos $a_1\in[7,9]$, $a_2\in[42,46]$; $x_1+x_2=100$; precios $0{,}22$ y $0{,}40$~€/kg.
\planteamiento para una restricción $\ge$ con $x\ge 0$, el peor caso usa las cotas inferiores (\cref{prop:robusta} con el signo cambiado).
\pasos
\begin{enumerate}[label=\arabic*.]
  \item (a) Contrapartida robusta: $7x_1+42x_2\ge 1800$.
  \item (b) Robusta: $7(100-x_2)+42x_2\ge 1800$, $35x_2\ge 1100$, $x_2\ge 31{,}43$~kg. Nominal ($a=(8,44)$): $800+36x_2\ge 1800$, $x_2\ge 27{,}78$~kg.
  \item (c) Coste nominal: $0{,}22\cdot 72{,}22+0{,}40\cdot 27{,}78=27{,}00$~€. Robusto: $0{,}22\cdot 68{,}57+0{,}40\cdot 31{,}43=27{,}66$~€.
\end{enumerate}
\resultado \textbf{soja mínima 31,43~kg (robusta) frente a 27,78~kg (nominal); precio de la robustez 0,66~€ por cada 100~kg (2,4\,\%).}
\comprobacion con la mezcla robusta y los contenidos mínimos: $7\cdot 68{,}57+42\cdot 31{,}43=480+1320=1800$.
\end{solucion}

\begin{solucion}{eje:6-saa}
\datos escenarios $64,72,88,95,110$ (ordenados), fractil $0{,}8214$.
\planteamiento con la distribución empírica, el \cref{teo:fractil} discreto da el menor escenario con $\hat F\ge 0{,}8214$.
\pasos
\begin{enumerate}[label=\arabic*.]
  \item $\hat F(95)=4/5=0{,}8<0{,}8214$; $\hat F(110)=1$: $q_{\mathrm{SAA}}=110$~hL.
  \item Normal ajustada: media $85{,}8$, desviación muestral $18{,}31$; $q=85{,}8+0{,}9208\cdot 18{,}31=102{,}7$~hL.
  \item Con cinco escenarios, el cuantil empírico salta de 95 a 110 y depende de un único valor extremo; la decisión tiene mucha varianza y su beneficio dentro de muestra está sesgado al alza.
\end{enumerate}
\resultado \textbf{SAA: 110~hL; normal ajustada: 102,7~hL}; con tan pocos escenarios, la decisión es inestable y hay que validarla fuera de muestra.
\comprobacion las transparencias de EE364a recomiendan aumentar el número de muestras hasta que el valor fuera de muestra se estabilice.
\end{solucion}

\begin{solucion}{eje:6-fractil}
\datos $D$ discreta con valores $d_1<\dots<d_K$ y probabilidades $p_k$.
\planteamiento estudiar $G(q)=\sum_kp_k[c_u(d_k-q)^{+}+c_o(q-d_k)^{+}]$ entre dos valores consecutivos.
\pasos
\begin{enumerate}[label=\arabic*.]
  \item Cada término es lineal en $q$ salvo en $d_k$, así que $G$ es lineal a trozos con vértices en los $d_k$.
  \item En $(d_k,d_{k+1})$, las demandas $\le d_k$ son de sobrante y las demás de falta: pendiente $c_oF(d_k)-c_u(1-F(d_k))=(c_u+c_o)F(d_k)-c_u$.
  \item La pendiente es no decreciente en $k$ (convexidad) y pasa de negativa a no negativa en el primer $k$ con $F(d_k)\ge c_u/(c_u+c_o)$.
  \item El mínimo de una función convexa lineal a trozos está en el vértice donde cambia el signo de la pendiente: $q^{*}=d_k$ para ese $k$.
\end{enumerate}
\resultado \textbf{$q^{*}=\min\{d_k: F(d_k)\ge c_u/(c_u+c_o)\}$.}
\comprobacion en el \cref{eje:6-vss}, $F(100)=0{,}7<0{,}8214\le F(140)$ da $q^{*}=140$, el valor obtenido allí.
\end{solucion}

\begin{solucion}{eje:6-maldicion}
\datos $\E\hat J_i=J_i$; en (b), $\hat J_i=J+\varepsilon_i$ con $\varepsilon_i\sim\mathcal N(0,\sigma^{2})$ independientes.
\planteamiento (a) desigualdad entre el máximo de esperanzas y la esperanza del máximo; (b) esperanza del máximo de dos normales.
\pasos
\begin{enumerate}[label=\arabic*.]
  \item (a) Para cada $j$, $\max_i\hat J_i\ge\hat J_j$; tomando esperanzas, $\E[\max_i\hat J_i]\ge J_j$ para todo $j$, luego $\E[\max_i\hat J_i]\ge\max_jJ_j$.
  \item (b) $\max(\varepsilon_1,\varepsilon_2)=\tfrac12(\varepsilon_1+\varepsilon_2)+\tfrac12\abs{\varepsilon_1-\varepsilon_2}$. Como $\varepsilon_1-\varepsilon_2\sim\mathcal N(0,2\sigma^{2})$ y $\E\abs{Z}=s\sqrt{2/\pi}$ para $Z\sim\mathcal N(0,s^{2})$: $\E\max=\tfrac12\sqrt{2}\sigma\sqrt{2/\pi}=\sigma/\sqrt{\pi}\approx 0{,}564\,\sigma$.
  \item (c) Separar los datos: elegir la alternativa con un conjunto de escenarios o semanas y estimar su valor con otro independiente; o usar estimadores con contracción hacia la media.
\end{enumerate}
\resultado \textbf{(a) $\E[\max\hat J]\ge\max J$; (b) sesgo $\sigma/\sqrt{\pi}\approx 0{,}564\,\sigma$ con dos alternativas iguales; (c) validar con datos no usados para elegir.}
\comprobacion una simulación con $10^{6}$ pares de normales estándar da una media del máximo de $0{,}565$, frente al valor teórico $1/\sqrt{\pi}=0{,}564$.
\end{solucion}
